% ============================================================================
% 2000 年春季高考（京、皖）文科 第 21 题 图二：折叠后三棱锥 D′-ABC（直二面角）
% ----------------------------------------------------------------------------
% 来源：.work/phase4-b20/source/2000春季京皖文.pdf
%       SHA-256 b9552b8354735296f6ea1ed051085e094e4af0c14ff97ada2dbd4e579b14287f
%       第 1 页，内容流对象 47（图注「图二」不在本图内）
% 重绘方式：qpdf --qdf 取内容流 → 解析 m/l/c/re/S/f、w、d → 本地坐标
%          本地坐标 = 页面坐标 − (970.740, 359.960)，原点取顶点 A；
%          1 单位 = 1 TeX pt；页面为 A3（1190.55 × 841.89 pt）。
% 元素计数：实线路径 3（AB、BD′、D′A），虚线路径 3（AC、BC、CD′），标签 4（A, B, C, D′）
% 线宽：全部 0.40 pt；虚线 on/off 各 2.99 pt
% 坐标基准与源位的偏差：端点四舍五入到 0.001 pt；全部为直线段，无简化。
% 标签：位置 = pdftotext -bbox 文本中心 − 基准；C、D′ 取「字母 + 撇」合成包围盒中心（撇号在内容流中为独立 F9 字形，pdftotext 未单列）；C 源位 (21.920,14.130)，右移 2.4 pt、下移 3.4 pt；其余未动。净距均 ≥ 2 pt
% ============================================================================
\begin{tikzpicture}[x=1pt, y=1pt,
    every node/.style={inner sep=0pt, font=\fontsize{10.5pt}{12.6pt}\selectfont}]
  % --- 实线：AB、BD′、D′A
  \draw[line width=0.4pt] (0,0) -- (68.040,0);
  \draw[line width=0.4pt] (68.040,0) -- (17.010,34.020);
  \draw[line width=0.4pt] (17.010,34.020) -- (0,0);
  % --- 虚线：AC、BC、CD′
  \draw[line width=0.4pt, dash pattern=on 2.99pt off 2.99pt] (0,0) -- (20.410,20.410);
  \draw[line width=0.4pt, dash pattern=on 2.99pt off 2.99pt] (68.040,0) -- (20.410,20.410);
  \draw[line width=0.4pt, dash pattern=on 2.99pt off 2.99pt] (20.410,20.410) -- (17.010,34.020);
  % --- 标签
  \node at (-4.930,-5.890) {$A$};
  \node at (72.710,-5.890) {$B$};
  \node at (16.760,41.350) {$D'$};
  \node at (24.320,10.730) {$C$};
  % 图注（源页图下居中；主会话添加）
  \node at (34.020,-12.000) {图二};
\end{tikzpicture}
% ----------------------------------------------------------------------------
% 自验：A=(0,0)、B=(68.040,0)、D′=(17.010,34.020)、C=(20.410,20.410)，与内容流 970.74/1038.78/987.75/991.15、359.96/393.98/380.37 逐点吻合；D′ 到 A、B 的距离分别 38.03、61.33 pt，C 由三条虚线汇于一点；标签：A、B、D′ 取源位；C 合成中心 (992.66,374.09) 右移 2.4 pt、下移 3.4 pt。净距 A 3.31、B 3.23、D′ 3.71、C 2.26 pt。探针编译：0 错误、0 警告。
% ============================================================================
